\documentclass{article}
\usepackage{listings}

\title{Contributing to KDEEcoLab - structure}

\begin{document}

% turn off pagenumbering
\pagenumbering{gobble}
\maketitle

\section{Get familiar with the project}
\subsection{Read all ressources (GitLab repository, Project's handbook)}
\subsection{Run the project locally and test the scripts}
\subsection{Change the source code and perform a merge request}
\subsection{Introduction to the project and topic}

\section{Selenium}
\subsection{Evaluation of the framework for fulfilling the task}
\subsection{Learning Selenium}
\subsubsection{How does it work}
\subsubsection{Building first test scripts}

\section{Transfer scripts to the target framework Selenium for each application and test-case}
\subsection{Baseline, Idle, Standard Usage Scenario}
\subsubsection{Okular}
\subsubsection{Kate}
\subsubsection{LibreOffice}

\section{Setup Testing Lab}
\subsection{Read documentation of the software energy consumption measurement process}
\subsection{Choose Software and Hardware components and plan the setup}
\subsection{Interfaces between used components}
\subsection{Test the Lab-setup}

\section{Testing software for it's energy consumption}
\subsection{Making actual testruns for each scenario}
\subsection{Process the collected data}

\section{Evaluation}
\subsection{Evaluation of the testing scripts}
\subsection{Evaluation of the measurement results}

\newpage

\begin{center}
\section*{Rough timeline for project points:}
\begin{tabular}{ | l | r | }
 \hline
 Task & Deadline \\
 \hline\hline
 1: Get familiar with the project and perform a successful merge request   & Nov \\
 \hline
 2: Evaluate Selenium and learn the important parts of the framework & Dec\\
 \hline
 3: Translation of Ocular, Kate and LibreOffice to Selenium & Feb\\
 \hline
 4: Setup Testing Lab & Mar\\
 \hline
 5: Test software's energy consumption & Apr\\
 \hline
 6: Evalutaion & May\\
 \hline
\end{tabular}
\end{center}

\begin{center}
\section*{Weekly Schedule}
\begin{tabular}{ | c | c | c | }
 \hline
 Week & Task & Meeting \\
 \hline\hline
 41  & Setup fedora VM (with KDE Plasma) and ydotool/kdotool & y \\
 \hline
 42  & Setup the project locally and and test the local setup/run and test scripts & y \\
 \hline
 43  & Change something and perform first merge request to the repository & y \\
 \hline
 44  & Write introduction to the topic (chapter 1) & \\
 \hline
 45  & Write introduction to the topic (chapter 1) & \\
 \hline
 46  & Learn Selenium, AT-SPI and needed components & y \\
 \hline
 47  & Setup Selenium, AT-SPI & \\
 \hline
 48  &  Setup Selenium, AT-SPI & y \\
 \hline
 49  &  Write and test first Selenium scripts & \\
 \hline
 50  &  Translation Okular to Selenium - baseline, idle & y \\
 \hline
 51  & Translation Okular to Selenium - sus & \\
 \hline
 52  & Translation Okular to Selenium - sus; Test the scripts, evaluation with team & \\
 \hline
 53  & buffer & \\
 \hline
 1  & Translation Kate to Selenium - baseline, idle & y \\
 \hline
 2  & Translation Kate to Selenium - sus & \\
 \hline
 3  & Translation Kate to Selenium - sus; Test the scripts, evaluation with team & y \\
 \hline
 4  & Write thesis & \\
 \hline
 5  & Write thesis & \\
 \hline
 6  & Plan Lab - create checklist & ? \\
 \hline
 7  & (Holidays - school closed): Plan Lab, write thesis &  \\
 \hline
 8  & Setup and Test Lab - troubleshooting &  \\
 \hline
 9  & Setup and Test Lab - troubleshooting & ? \\
 \hline
 10  & Perform tests for each of the three scenarios & \\
 \hline
 11  & Perform tests for each of the three scenarios & ? \\
 \hline
 12  & ?Perform tests?/Process test data & \\
 \hline
 13 & Evaluate tests & ? \\
 \hline
 14 & Write thesis & \\
 \hline
 15 & Eventually: scripts for new applications??? & \\
 \hline
 tba & tba & tba \\
 \hline

\end{tabular}
\end{center}

\begin{lstlisting}[language=Python]
import numpy as np

def incmatrix(genl1,genl2):
    m = len(genl1)
    n = len(genl2)
    M = None #to become the incidence matrix
    VT = np.zeros((n*m,1), int)  #dummy variable

    #compute the bitwise xor matrix
    M1 = bitxormatrix(genl1)
    M2 = np.triu(bitxormatrix(genl2),1)

    for i in range(m-1):
        for j in range(i+1, m):
            [r,c] = np.where(M2 == M1[i,j])
            for k in range(len(r)):
                VT[(i)*n + r[k]] = 1;
                VT[(i)*n + c[k]] = 1;
                VT[(j)*n + r[k]] = 1;
                VT[(j)*n + c[k]] = 1;

                if M is None:
                    M = np.copy(VT)
                else:
                    M = np.concatenate((M, VT), 1)

                VT = np.zeros((n*m,1), int)

    return M
\end{lstlisting}

\lstinputlisting[language=Python]{main.py}

\end{document}
